\paragraph{Idee und Erwartung.} Der Nutzer schlug vor: \emph{„Es kann auch helfen, bestimmte Bereiche zu schwächen, so
dass sich die Last gleichmäßiger verteilt.“} Das Prinzip ist aus dem Leichtbau bekannt: Kraft fließt dem steifsten
Weg nach; wer einen überlasteten Bereich nachgiebiger macht, schickt Last in unterbeanspruchte Bereiche. Da die
Spitzen an den Verbinderenden sitzen, wurde als einfachste Form geprüft, alle Verbinder von Entwurf 0169 weicher zu
machen (\SI{140}{}, \SI{100}{}, \SI{50}{GPa} statt \SI{183}{GPa}), jeweils im Standardnetz und mit 1,5-mal mehr
Stationen. Erwartet war (F-139): Spitze $-30$ bis $-45$~\%, $\sigma_{99}$ $-5$ bis $-15$~\%, Verformung $+20$ bis
$+60$~\%.

\begin{table}[H]\centering\scriptsize\begin{tabular}{@{}lrrrrrrrl@{}}\toprule
Nr. & E Verbinder & Stationen & $\sigma_{99}$ & max. & Verformung & Dehnung & $B\cdot n$ & alle Grenzen \\
 & & & MPa & MPa & mm & \% & ‰ & \\\midrule
0105 & 50 GPa & 64 & \num{358} & \num{743} & \num{27} & \num{0.17} & \num{2.13} & ja \\
0110 & 100 GPa & 64 & \num{331} & \num{734} & \num{22} & \num{0.16} & \num{2.13} & ja \\
0114 & 140 GPa & 64 & \num{318} & \num{725} & \num{21} & \num{0.15} & \num{2.13} & ja \\
0205 & 50 GPa & 96 & \num{372} & \num{1541} & \num{29} & \num{0.19} & \num{2.13} & ja \\
0210 & 100 GPa & 96 & \num{343} & \num{1398} & \num{25} & \num{0.18} & \num{2.13} & ja \\
0214 & 140 GPa & 96 & \num{330} & \num{1323} & \num{23} & \num{0.17} & \num{2.13} & ja \\
\bottomrule\end{tabular}
\caption{0169 mit weicheren Verbindern. Nummern:
1xx Standardnetz, 2xx 1,5-mal mehr Stationen; Original \SI{183}{GPa}: \SI{308}{MPa} bzw.\ \SI{319}{MPa}.}\end{table}

\paragraph{Ergebnis.} Die Spannung im Gehäuse steigt mit sinkender Verbindersteifigkeit gleichmäßig an: im
Standardnetz von \SI{308}{MPa} (\SI{183}{GPa}) über \SI{318}{MPa} und \SI{331}{MPa} auf \SI{358}{MPa}
(\SI{50}{GPa}, $+16$~\%); im feinen Netz von \SI{319}{MPa} auf \SI{372}{MPa}. Auch der Höchstwert steigt. Die
Verformung wächst von \SI{19}{mm} auf \SI{27}{mm} ($+41$~\%, im erwarteten Bereich), die Dehnung im Wickelpaket von
0,15~\% auf 0,17~\%. Entlastet werden nur die Verbinder selbst (99-\%-Wert \SI{192}{\to}\SI{127}{MPa}). Alle
Varianten bleiben zulässig, aber mit weniger Reserve.

\paragraph{Deutung.} Die Erwartung ist widerlegt, und der Grund ist lehrreich: Die Verbinder sind in diesem Tragwerk
nicht die Ursache der Last, sondern der Entlastungsweg. Sie verkürzen die freie Biegelänge der Gehäuse. Werden sie
insgesamt weicher, biegen sich die Gehäuse wieder stärker. Die Spitze am Verbinderende ist dagegen eine Folge der
\emph{Form} des Anschlusses (scharfe Kante, Steifigkeitssprung). Die Nutzer-Idee bleibt richtig, wenn sie örtlich
angewendet wird: ein Verbinder mit weichen Enden und steifem Kern -- wie eine auslaufende Klebefuge oder ein
angeschrägter Steg -- leitet die Kraft allmählich ein, ohne den Entlastungsweg aufzugeben. Dafür muss jeder Verbinder
in mehrere Körper mit eigener Steifigkeit zerlegt werden; die Steifigkeiten lassen sich dann nach dem
Spannungsverhältnis iterieren (fully stressed design).
